%======================================================================
% فصل دوم: مروری بر پیشینه
%======================================================================
\chapter{مروری بر پیشینه}
\label{ch:related}

\section{مقدمه}

در این فصل، پیشینه‌ی نظری و عملی مسئله‌ی یکپارچه‌سازی سرویس‌های هوش
مصنوعی بررسی می‌شود. مسئله‌ی اصلی این پروژه صرفاً دسترسی به یک مدل
هوشمند نیست، بلکه مدیریت هم‌زمان چند ارائه‌دهنده، یکسان‌سازی رابط‌های
برنامه‌نویسی، کنترل دسترسی، صورت‌حساب‌گذاری و گزارش‌گیری در یک بستر واحد
است. به همین دلیل، مرور پیشینه باید هم مفاهیم معماری نرم‌افزار و هم
نمونه‌سامانه‌های عملی موجود در بازار را در بر گیرد.

در سال‌های اخیر، \term{رابط‌های برنامه‌نویسی کاربردی}{Application Programming
Interface (API)} به روش غالب دسترسی به مدل‌های هوش مصنوعی تبدیل شده‌اند.
از سوی دیگر، رشد سریع مدل‌های زبانی بزرگ باعث شده است که ارائه‌دهندگان
متعدد، هر یک با قرارداد ارتباطی، قالب داده، سیاست مالی و محدودیت‌های
اجرایی خاص خود، وارد این حوزه شوند. نتیجه آن است که مسئله‌ی
\term{یکپارچه‌سازی}{Integration} اکنون به اندازه‌ی خود مدل‌ها اهمیت دارد.

\section{مفاهیم پایه}

\subsection{سامانه‌ی واسط \lr{API}}

\term{سامانه‌ی واسط}{Gateway} در معماری نرم‌افزار، لایه‌ای است که بین
مصرف‌کنندگان سرویس و سامانه‌های پشتیبان قرار می‌گیرد و مسئولیت‌هایی مانند
احراز هویت، اعتبارسنجی درخواست، محدودسازی نرخ، راهبری، ثبت گزارش و
یکسان‌سازی پاسخ را بر عهده می‌گیرد. در معماری‌های وب مبتنی بر
\term{سبک انتقال حالت بازنمایی}{Representational State Transfer (REST)}،
این لایه نقش مهمی در کاهش وابستگی مشتری به جزئیات داخلی سرویس‌های پشت
صحنه دارد~\cite{fielding2000}.

در مسئله‌ی حاضر، سامانه‌ی واسط به جای آن‌که تنها یک \lr{Reverse Proxy}
باشد، یک لایه‌ی منطقی کامل است. این لایه باید تفاوت میان قالب‌های
درخواست و پاسخ ارائه‌دهندگان مختلف را پنهان کند، دسترسی کاربران را
مدیریت نماید، وضعیت مالی را کنترل کند و گزارش‌های قابل استفاده برای
کاربر و مدیر سامانه تولید نماید. بنابراین، سامانه‌ی واسط مورد نظر این
پروژه ترکیبی از وظایف \lr{API Gateway}، \lr{Billing Gateway} و
\lr{Integration Layer} است.

\subsection{معماری خدمات‌گرا و ریزسرویس}

در معماری‌های خدمات‌گرا و ریزسرویسی، تفکیک مسئولیت‌ها و مرزبندی روشن
اجزای سامانه، یکی از اصول کلیدی برای مقیاس‌پذیری و نگهداشت‌پذیری است.
به‌جای آن‌که منطق احراز هویت، مسیردهی، صورت‌حساب و گزارش‌گیری در یک لایه‌ی
یکپارچه و فشرده پیاده‌سازی شوند، هر حوزه به‌صورت ماژول مستقل طراحی
می‌گردد~\cite{newman2015,richardson2018}. مزیت چنین رویکردی آن است که
تغییر در یک حوزه اثر جانبی کمتری بر سایر حوزه‌ها می‌گذارد، آزمون‌پذیری
سامانه افزایش می‌یابد زیرا هر بخش مرز رفتاری روشن‌تری پیدا می‌کند، و
افزودن قابلیت‌های جدیدی مانند راهبری هوشمند یا مدار قطعه بدون بازنویسی
ساختار اصلی امکان‌پذیر می‌شود.

هرچند پروژه‌ی حاضر در قالب یک سامانه‌ی متمرکز \lr{Laravel} پیاده‌سازی شده
است، اما در طراحی داخلی آن از همین اصل تفکیک مسئولیت استفاده شده و
لایه‌های دامنه، سرویس، مخزن، منبع و خط لوله از یکدیگر جدا شده‌اند. این
موضوع باعث می‌شود ساختار سامانه، از منظر معماری، رفتاری نزدیک به یک
معماری خدمات‌گرا داشته باشد.

\subsection{مدل‌های زبانی بزرگ و سرویس‌های هوش مصنوعی}

تحول اصلی در سرویس‌های هوش مصنوعی معاصر، با فراگیر شدن معماری
\term{ترنسفورمر}{Transformer} و مدل‌های زبانی بزرگ رخ داد~\cite{vaswani2017}.
این مدل‌ها به‌واسطه‌ی آموزش روی داده‌های وسیع و پشتیبانی از الگوهای
\term{چندوجهی}{Multimodal}، امروز در کاربردهایی مانند گفتگو، خلاصه‌سازی،
جاسازی متن، تولید تصویر و تبدیل گفتار به متن به‌کار می‌روند.

با وجود شباهت کارکردی، سرویس‌های مبتنی بر این مدل‌ها از نظر قرارداد
فنی تفاوت‌های محسوسی دارند. برای نمونه، \lr{OpenAI} مجموعه‌ای از نقاط
پایانی نسبتاً استاندارد برای گفتگو، \lr{responses}، جاسازی متن، تصویر و
صوت ارائه می‌دهد~\cite{openai2024}. در مقابل، \lr{Gemini} از \lr{API}
بومی خود استفاده می‌کند و نقاط پایانی اصلی آن \lr{generateContent} و
\lr{embedContent} هستند~\cite{gemini2024}. در نتیجه، قالب پیام‌ها و
فراساختار مصرف در این سرویس متفاوت است. این تفاوت‌ها نشان می‌دهند که
یک سامانه‌ی واسطِ مؤثر باید هم از
\term{سازگاری بیرونی}{External Compatibility} و هم از
\term{نگاشت درونی}{Internal Mapping} پشتیبانی کند.

\subsection{رابط‌های سازگار با \lr{OpenAI}}

در عمل، رابط \lr{OpenAI} به یکی از مراجع غیررسمی این حوزه تبدیل شده است.
بسیاری از ابزارهای مشتری، کتابخانه‌ها و حتی برخی ارائه‌دهندگان دیگر، برای
کاهش هزینه‌ی مهاجرت مشتریان، خود را با این قرارداد بیرونی سازگار
می‌کنند~\cite{openai2024,groq2024,openrouter2024}. چنین روندی از یک سو
هزینه‌ی تغییر ارائه‌دهنده را کاهش می‌دهد، اما از سوی دیگر تفاوت در
قابلیت‌ها، مدل‌ها و فیلدهای مصرف را از بین نمی‌برد؛ بنابراین این
ناهمگنی باید در لایه‌ی واسط مدیریت شود.

بر همین اساس، استفاده از یک رابط بیرونی سازگار با \lr{OpenAI} در این
پروژه یک تصمیم معماری راهبردی بوده است؛ زیرا هم برای توسعه‌دهندگان
آشناست و هم امکان پشتیبانی از چند ارائه‌دهنده را با حداقل تغییر در
سمت مشتری فراهم می‌کند.

\section{سیستم‌های مشابه}

\subsection{\lr{OpenAI}}

\lr{OpenAI} یکی از مهم‌ترین بازیگران این حوزه است و مستندات آن عملاً
به مرجع طراحی بسیاری از مشتریان و سامانه‌های واسط تبدیل شده‌اند
~\cite{openai2024}. مزیت اصلی این بستر، انسجام رابط برنامه‌نویسی،
پوشش چند نوع خدمت هوش مصنوعی و فراگیری ابزارهای پیرامونی آن است.

با این حال، \lr{OpenAI} ذاتاً یک ارائه‌دهنده‌ی منفرد است، نه یک لایه‌ی
تجمیع‌کننده‌ی چند ارائه‌دهنده. بنابراین مسائلی مانند مهاجرت میان
ارائه‌دهندگان، مقایسه‌ی قیمت‌ها، اعمال سیاست‌های مالی داخلی سازمان و
مدیریت محلی کاربران، خارج از دامنه‌ی اصلی آن قرار دارند.

\subsection{\lr{Gemini}}

\lr{Gemini} به‌عنوان بستر بومی \lr{Google DeepMind}، قابلیت‌های زبانی و
چندوجهی قدرتمندی ارائه می‌کند، اما از نظر شکل \lr{API} با الگوی
\lr{OpenAI} یکسان نیست~\cite{gemini2024}. درخواست‌ها در این بستر معمولاً
بر پایه‌ی ساختار \lr{contents/parts} و پاسخ‌ها با قالبی متفاوت از
فیلدهای متعارف \lr{OpenAI} ارائه می‌شوند.

از منظر این پایان‌نامه، اهمیت \lr{Gemini} در آن است که نشان می‌دهد
یکپارچه‌سازی چند ارائه‌دهنده فقط مسئله‌ی تغییر \lr{Base URL} نیست؛ بلکه
در برخی موارد به نگاشت ساختاری درخواست، پاسخ و متاداده‌ی مصرف نیاز دارد.
بنابراین، پشتیبانی از \lr{Gemini} محک خوبی برای ارزیابی انعطاف‌پذیری
معماری سامانه‌ی واسط است.

\subsection{\lr{Groq}}

\lr{Groq} بستری است که با تأکید بر سرعت بالای استنتاج و با ارائه‌ی
\lr{API} سازگار با \lr{OpenAI} شناخته می‌شود~\cite{groq2024}. مزیت اصلی آن
برای توسعه‌دهنده این است که مهاجرت از مشتریان مبتنی بر \lr{OpenAI}
به‌سادگی انجام می‌شود و در بسیاری از موارد، تغییر اصلی فقط در
\code{provider}، \code{model} و اطلاعات احراز هویت خلاصه می‌شود.

با وجود این سازگاری، \lr{Groq} نیز یک تجمیع‌کننده‌ی چندارائه‌دهنده نیست و
منطق مدیریت کاربران، کیف پول، اشتراک و گزارش‌های سازمانی را بر عهده
ندارد. از این رو، شباهت آن به پروژه‌ی حاضر بیشتر در لایه‌ی قرارداد
فنی \lr{API} است تا در حوزه‌ی مدیریت سرویس.

\subsection{\lr{OpenRouter}}

\lr{OpenRouter} از نزدیک‌ترین نمونه‌های تجاری به ایده‌ی این پروژه است؛
زیرا یک رابط نسبتاً واحد برای دسترسی به چندین ارائه‌دهنده و مدل مختلف
فراهم می‌کند~\cite{openrouter2024}. این بستر نشان می‌دهد که مفهوم
\term{تجمیع هوش مصنوعی}{AI Aggregation} در عمل نیز دارای تقاضای واقعی است.

با این حال، تمرکز اصلی \lr{OpenRouter} بر عرضه‌ی یک \lr{API} عمومی برای
توسعه‌دهندگان است، نه بر مدیریت داخلی یک سازمان. در نتیجه، ویژگی‌هایی
مانند احراز هویت مبتنی بر \lr{OTP}، کیف پول کاربرمحور، نقش‌ها و سطوح دسترسی،
ثبت ممیزی عملیاتی و گزارش‌های مالی سفارشی، در سطحی که برای یک سامانه‌ی
سازمانی لازم است، در مرکز طراحی آن قرار ندارند.

\subsection{\lr{Anyscale Endpoints}}

\lr{Anyscale Endpoints} نیز بستری برای ارائه و میزبانی مدل‌ها در قالب
خدمات قابل فراخوانی است~\cite{anyscale2024}. نقطه‌ی قوت این بستر بیشتر در
حوزه‌ی استقرار، مقیاس‌پذیری عملیاتی و خدمت‌رسانی مدل‌ها است. در نتیجه،
تمرکز آن با سامانه‌ی حاضر تفاوت دارد: \lr{Anyscale} بیشتر مسئله‌ی
\term{خدمت‌رسانی مدل}{Model Serving} را حل می‌کند، در حالی که پروژه‌ی
حاضر مسئله‌ی \term{یکپارچه‌سازی دسترسی و مدیریت مصرف} را هدف قرار داده است.

\section{شکاف پژوهشی و معیارهای مقایسه}

بررسی پیشینه نشان می‌دهد که در این مسئله با دو دسته منبع روبه‌رو هستیم.
دسته‌ی نخست، منابع معماری نرم‌افزار هستند که اصولی مانند \lr{REST}،
تفکیک مسئولیت، و الگوهای خدمات‌گرا را توضیح می‌دهند
~\cite{fielding2000,newman2015,richardson2018}. دسته‌ی دوم، مستندات رسمی
ارائه‌دهندگان هستند که قراردادهای واقعی \lr{API}، مدل‌های قیمت‌گذاری و
قابلیت‌های اجرایی را مشخص می‌کنند
~\cite{openai2024,gemini2024,groq2024,openrouter2024,anyscale2024}. از آن‌جا
که موضوع این پایان‌نامه یک مسئله‌ی پیاده‌سازی‌محور و وابسته به قراردادهای
واقعی سرویس‌ها است، استفاده از مستندات رسمی در کنار منابع معماری یک انتخاب
روش‌مند است، نه صرفاً یک اتکای توصیفی.

برای آن‌که مقایسه‌ی سامانه‌های مشابه صرفاً فهرست‌وار نباشد، در این
پایان‌نامه پنج معیار تحلیلی برای سنجش آن‌ها در نظر گرفته می‌شود. این
معیارها عبارت‌اند از سازگاری بیرونی، یعنی میزان نزدیکی رابط سامانه به
قرارداد \lr{OpenAI} و هزینه‌ی مهاجرت مشتریان؛ پوشش چندارائه‌دهنده، یعنی
این‌که سامانه واقعاً چند ارائه‌دهنده را تجمیع می‌کند یا صرفاً بستری برای
یک ارائه‌دهنده است؛ عمق نگاشت درونی، یعنی توانایی سامانه در پنهان‌سازی
تفاوت‌های ساختاری میان ارائه‌دهندگان سازگار و ناسازگار؛ پشتیبانی مدیریتی
و مالی، یعنی وجود قابلیت‌هایی مانند احراز هویت، مدیریت کلید، کنترل
دسترسی، کیف پول، اشتراک و گزارش‌گیری؛ و در نهایت جهت‌گیری استقرار، یعنی
این‌که سامانه برای مصرف عمومی توسعه‌دهنده، میزبانی مدل یا مدیریت مصرف در
سطح سازمان طراحی شده است.

بر اساس این معیارها، شکاف اصلی پیشینه را می‌توان چنین صورت‌بندی کرد:
\textbf{سامانه‌های موجود یا روی ارائه‌ی مستقیم مدل تمرکز دارند، یا روی
تجمیع عمومی دسترسی؛ اما ترکیب هم‌زمان قرارداد بیرونی آشنا، نگاشت چند
ارائه‌دهنده، و کنترل مالی و مدیریتی در سطح سازمان کمتر دیده می‌شود.}
این شکاف، مبنای طراحی سامانه‌ی پیشنهادی در فصل بعدی است.

\section{مقایسه و جمع‌بندی}

مرور سامانه‌های فوق نشان می‌دهد که هر یک بخشی از مسئله را پوشش می‌دهند،
اما هیچ‌کدام دقیقاً مجموعه‌ی نیازهای این پروژه را به‌صورت یکجا پاسخ
نمی‌دهند. جدول \ref{tab:related-comparison} خلاصه‌ای از این مقایسه را
نشان می‌دهد.

\begin{table}[H]
\centering
\footnotesize
\caption{مقایسه‌ی سامانه‌های بررسی‌شده با مسئله‌ی این پایان‌نامه}
\label{tab:related-comparison}
\renewcommand{\arraystretch}{1.4}
\begin{tabularx}{\textwidth}{>{\raggedleft\arraybackslash}p{2.2cm} >{\centering\arraybackslash}p{1.9cm} >{\centering\arraybackslash}p{2cm} >{\centering\arraybackslash}p{2.1cm} X}
\toprule
\textbf{سامانه} & \textbf{\shortstack{چند\\ارائه‌دهنده}} & \textbf{\shortstack{سازگار با\\ \lr{OpenAI}}} & \textbf{\shortstack{مدیریت\\داخلی}} & \textbf{تمرکز اصلی} \\
\midrule
\lr{OpenAI} & خیر & بله & خیر & ارائه‌ی مستقیم سرویس‌های هوش مصنوعی \\
\lr{Gemini} & خیر & خیر & خیر & ارائه‌ی بومی مدل‌های چندوجهی \\
\lr{Groq} & خیر & بله & خیر & استنتاج سریع با قرارداد سازگار \\
\lr{OpenRouter} & بله & بله & نسبی & تجمیع عمومی چند مدل برای توسعه‌دهندگان \\
\lr{Anyscale Endpoints} & نسبی & نسبی & خیر & میزبانی و خدمت‌رسانی مدل‌ها \\
\textbf{سامانه‌ی این پایان‌نامه} & بله & بله & بله & یکپارچه‌سازی فنی، مالی و مدیریتی در سطح سازمان \\
\bottomrule
\end{tabularx}
\normalsize
\end{table}

در نتیجه، شکاف اصلی در پیشینه‌ی موجود را می‌توان چنین خلاصه کرد: سامانه‌ای
مورد نیاز است که از یک سو برای توسعه‌دهنده، رابطی آشنا و سازگار با
\lr{OpenAI} ارائه دهد و از سوی دیگر، نیازهای مدیریتی و مالی یک سازمان
را نیز پوشش دهد. این شکاف، مبنای طراحی سامانه‌ی پیشنهادی این پایان‌نامه
است که در فصل بعدی معماری و مدل داده‌ی آن تشریح می‌شود.
